The linear representation hypothesis---that concepts are encoded as directions in
the residual stream of large language models---underpins practical techniques such
as activation steering and concept editing. Practitioners routinely compose
steering vectors via addition (\eg adding ``formal'' and ``concise'' directions),
yet no systematic study has established \emph{when} such composition succeeds.
We conduct the first controlled empirical study of linear-direction
compositionality, testing 10 concept pairs across 5 semantic categories
(hierarchical, independent, semantic, syntactic, and entangled) in
\pythia using three complementary metrics: pairwise orthogonality,
probing-based information preservation, and steering-based component
preservation. We find that most concept directions extracted via
mean-difference are surprisingly non-orthogonal (median $\abscossim = 0.91$ at
layer~20), yet composition quality varies dramatically by metric.
Probing-based preservation is uniformly high (${\sim}0.95$), masking sharp
differences revealed by steering: near-orthogonal pairs ($\abscossim < 0.3$)
preserve both components well (mean~0.68), while highly aligned pairs
($\abscossim > 0.9$) often lose one or both (mean~0.55). The absolute cosine
similarity between directions is the strongest predictor of steering
composition success, more reliable than the semantic category of the concept
pair. Our results both support and qualify the linear representation
hypothesis: linear directions exist and are meaningful, but they do not form
a clean vector space amenable to free composition.
